\chapter{Literature Review}
\label{chap:literature}

\section{Scope of the Review}

This chapter reviews work relevant to \textbf{Flowmatic}: assessing traffic data quality, automating preparation, deploying stream-oriented ITS platforms, and evaluating neural models for forecasting, reconstruction, imputation, and classification.

\section{Data Quality in Transportation Analytics}

General data quality research treats quality as fitness for use~\cite{wang1996beyond,pipino2002data}. Transportation streams are spatially distributed, time ordered, and heterogeneous~\cite{rahm2000data,heinrich2007assessing}. Composite indices~\cite{bieberstein2006data} aggregate dimensions when weights remain application dependent. Flowmatic implements six dimensions---completeness, consistency, accuracy, timeliness, uniqueness, and validity---in \texttt{QualityService} (Table~\ref{tab:dqi_weights}).

\section{Automated Preprocessing and Missing Data}

Automated preprocessing adapts operations to data and downstream tasks~\cite{liu2022automated,sharma2022scenario,feurer2019auto}. Missing traffic data requires temporal and spatial structure~\cite{tang2024missing,chan2023missing}. Phase~2 bundles in this thesis materialise windows for neural training; SAITS handles masked reconstruction within the production portfolio.

\section{Feature Engineering and Representation Learning}

Traffic analytics benefit from cyclical time features, rolling statistics, and graph structure~\cite{cai2022feature,luo2024spatiotemporal,wu2021graph}. Flowmatic generates windowed tensors and adjacency for STGCN while keeping schema-level features auditable in batch uploads.

\section{Anomaly Detection and Classification}

Classical outlier methods (LOF, Isolation Forest) and deep reconstruction models address different anomaly types~\cite{breunig2000lof,liu2008isolation,pang2021deep,ding2024variational}. Flowmatic deploys \textbf{TranAD-style} reconstruction on Astana streams (Equation~\ref{eq:tranad_score}) and a \textbf{Transformer} severity classifier with macro-F1 reporting (Equation~\ref{eq:macro_f1}).

\section{Intelligent Transportation Systems and Deployment}

Smart-city architectures combine sensing, streaming, and analytics~\cite{zanella2014iot,its_bigdata2018,event_driven2025}. Flowmatic uses HTTP/WebSocket ingestion, queue workers, and a model-inference microservice rather than a notebook-only workflow.

\section{Explainability}

Agencies require traceable diagnostics~\cite{lundberg2017unified,lipton2018mythos}. Flowmatic emphasises deterministic Insight Engine metrics plus neural residuals; optional LLM narration is constrained to pipeline context.

\section{Neural Time-Series Models}

Decomposition linear models~\cite{zeng2023transformers}, PatchTST~\cite{nie2023patchtst}, and STGCN~\cite{yu2018stgcn} motivate the eight-checkpoint portfolio evaluated under a single temporal split protocol (Chapter~\ref{chap:results}).

\section{Scope of This Thesis}

The dissertation reports completed work only:
\begin{enumerate}
  \item \textbf{Flowmatic platform} --- batch preparation, smart-city workbench, registry routing, Docker deployment;
  \item \textbf{Neural portfolio} --- eight Hugging Face checkpoints with multi-seed held-out test metrics and streaming latency benchmarks.
\end{enumerate}
All quantitative claims in Chapters~\ref{chap:results}--\ref{chap:conclusion} refer to artifacts in this repository or published Hub weights.
